\begin{document}
\maketitle


\begin{abstract}
This report introduces AutoIntent, a novel AutoML framework tailored for text classification tasks. Unlike existing solutions, AutoIntent offers end-to-end automation, encompassing embedding model selection, classifier optimization, and decision threshold tuning, all within a modular, sklearn-like interface. The framework is designed to support multi-label classification and out-of-scope detection, with built-in augmentation tools leveraging large language models. Through extensive experimentation across multiple benchmarks, AutoIntent demonstrates competitive performance. Its support for various transformer backbones, lightweight models, and configurable presets enables users to balance accuracy and resource consumption effectively. This work positions AutoIntent as a promising practical solution for NLP applications.
\end{abstract}


% keywords can be removed
\keywords{automated machine learning \and text classifiction \and text augmentation}


\section{Introduction}

Text classification remains a fundamental task in natural language processing, with applications ranging from intent detection in conversational systems \cite{weld2021surveyjointintentdetection} to content categorization and sentiment analysis \cite{TAHA2024100664}. Modern NLP has been revolutionized by transformer-based embedding models \cite{xiao2024cpackpackedresourcesgeneral}, which provide rich contextual representations of text. However, effectively utilizing these models for classification tasks requires careful consideration of multiple components: the choice of embedding model, the selection of appropriate classification algorithms, and the optimization of their hyperparameters.

Traditional machine learning approaches often require manual tuning of these components, which can be time-consuming and requires significant expertise. While automated machine learning (AutoML) frameworks \cite{baratchi2024automated} have emerged to automate this process, existing solutions for NLP tasks often lack comprehensive support for the full spectrum of hyperparameter optimization (for instance, choosing best embedding model) and tuning confidence thresholds for handling multi-label classification and out-of-scope (OOS) detection \cite{clinc150}.

This technical report introduces AutoIntent \footnote{\url{https://github.com/DeepPavlov/AutoIntent}}, a novel AutoML framework specifically designed for text classification tasks. AutoIntent addresses these limitations by providing an end-to-end pipeline that automatically optimizes all aspects of the classification process, from embedding model selection to machine learning model choice and threshold tuning. The framework offers a sklearn-like interface \cite{pedregosa2018scikitlearnmachinelearningpython} for ease of use. AutoIntent supports features such as multi-label classification and OOS detection, bridging the gap between AutoML and conversation systems.

\section{Background}

Automated machine learning, by definition, is a tool for automating routines like data splitting for validation, hyperparameter tuning, model ensembling, and model selection. The AutoML is highly relevant in the scenarios when machine learning tasks need to be solved by non-experts and in conjunction with no-code ML which is sometimes called <<ML as a service>> \cite{bisong2019overview,barga2015predictive,liberty2020elastic,ledell2020h2o,carney2020teachable,acito2023predictive}. The applications include sentiment analysis, robotics, healthcare, business analysis \cite{yuan2024automated,salehin2024automl}.

Tabular AutoML focuses on feature engineering, feature selection and model ensembling \cite{feurer2022autosklearn20handsfreeautoml, vakhrushev2022lightautomlautomlsolutionlarge, erickson2020autogluontabularrobustaccurateautoml, ledell2020h2o, nikitin2021automated}. Usually, they employ classical machine learning methods like GBMs \cite{chen2016xgboost,prokhorenkova2018catboost,ke2017lightgbm} and linear models, fast hyperparameter tuning methods with budget-aware strategies \cite{akiba2019optuna}, and ensembling strategies like stacking, blending and voting. Such frameworks sometimes can supersede exploratory data analysis and extensive research with just a running preset training recipe. It is not rare to see AutoML frameworks winning machine learning contests, but the open source solutions often are not transferable to production-ready systems as the resulting pipeline is an ensemble of a numerous amount of models without clear guides for deployment.

Neural architecture search can be viewed as automation in the field of deep learning \cite{salehin2024automl}. It emphasizes finding optimal computational graph using approaches like cell-based and hierarchical search spaces \cite{zoph2018learningtransferablearchitecturesscalable,real2019regularizedevolutionimageclassifier} or using scaling laws \cite{tan2020efficientnetrethinkingmodelscaling}. It cannot be treated as full-fledged AutoML, as it designed to address only the model selection problem. Though, it can be a part of an AutoML pipeline \cite{autokeras}.

NLP AutoML primarily stands out from other AutoML by native support of text inputs \cite{tang2024autogluonmultimodalautommsuperchargingmultimodal,vakhrushev2022lightautomlautomlsolutionlarge}. This is extra important for use by non-experts, as it removes the requirement of manual tokenization and vectorization. Though some frameworks provide auxiliary tools for text feature extraction \cite{ledell2020h2o}. The next peculiarity about text AutoML frameworks is their usage of transformer-based backbones, which is explainable, as this is the state-of-the-art in the field of NLP. Note that NLP AutoML primarily focuses on simple tasks like classification and regression, ignoring text generation and named entity recognition, for instance.

In AutoML frameworks, the model selection can be implemented with three ways. The first and most straightforward is to use hyperparameter tuning tools like optuna \cite{akiba2019optuna} or genetic algorithms \cite{feurer2022autosklearn20handsfreeautoml} with preset search spaces. Usually, these presets differ in how much time and computational resources they require to reach acceptable quality. The variety of presets is provided to cover all possible use cases and hardware settings. Another option is not to tune hyperparameters but use some generic hyperparameters that reach balance between the final quality and the generalization across different tasks. These hyperparameters can be obtained empirically \cite{tang2024autogluonmultimodalautommsuperchargingmultimodal}. The compromise between freezing hyperparameters and tuning all them is the meta-learning \cite{Desai2022,feurer2022autosklearn20handsfreeautoml,wang2021meta,tian2022meta,huisman2021survey}. In such approaches, there is a metamodel that takes a dataset as input and predicts hyperparameters for a machine learning model.

\section{AutoIntent}

\subsection{Design Principles}

The design of AutoIntent is guided by several key principles that address the limitations of existing NLP AutoML frameworks \ref{tab:frameworks} ensuring practical usability and maintainability.

\begin{enumerate}
    \item \textbf{Modular Architecture}: The framework is designed with a clear separation of concerns, where each component can function both independently and as part of the AutoML pipeline. This modularity facilitates:
    \begin{itemize}
        \item Independent development and testing of components
        \item Flexible deployment configurations
        \item Simplified maintenance and updates
        \item Clear and transparent framework structure
    \end{itemize}

    \item \textbf{Model Diversity}: To accommodate varying hardware constraints and performance requirements, the framework incorporates:
    \begin{itemize}
        \item Modern deep learning models for high-performance scenarios
        \item Classical machine learning models operating on transformer embeddings for resource-constrained environments
        \item A unified interface for all model types, ensuring consistent usage patterns
    \end{itemize}

    \item \textbf{Software Engineering Best Practices}: The framework adheres to established software development standards:
    \begin{itemize}
        \item Comprehensive documentation with user guides and API references
        \item Static type checking for improved code reliability
        \item Automated unit testing for robust quality assurance
        \item Unified code style enforced through linters and formatters
    \end{itemize}

    \item \textbf{Flexible Optimization Strategies}: The framework provides multiple approaches to model optimization:
    \begin{itemize}
        \item Hyperparameter optimization (HPO) presets for different computational budgets
        \item Non-HPO presets with generalizing HP configurations 
        \item Customizable search spaces for fine-grained control
    \end{itemize}

    \item \textbf{Conversational System Support}: The framework includes specialized features for dialog systems:
    \begin{itemize}
        \item Native support for multi-label classification
        \item Built-in out-of-scope (OOS) detection capabilities
        \item Few-shot learning optimizations
    \end{itemize}

    \item \textbf{Model Repository Integration}: The framework integrates HuggingFace model and datasets repository:
    \begin{itemize}
        \item Easy adoption of either the last SOTA pre-trained transformers or custom fine-tuned ones
        \item Version control for models
    \end{itemize}

    \item \textbf{Embedding-centric Design}: The framework adopts a transformer-based embedding approach:
    \begin{itemize}
        \item Elimination of traditional feature engineering stages as current embeddings are quite informative across variety of tasks \cite{enevoldsen2025mmtebmassivemultilingualtext}
        \item Utilization of pre-trained transformer embeddings in zero-shot manner
        \item Ability to select most appropriate source of features from HuggingFace repository
    \end{itemize}

    \item \textbf{Task Specialization}: The framework maintains a focused scope:
    \begin{itemize}
        \item Primary focus on intent classification
        \item Exclusion of unrelated tasks (e.g., computer vision, time series)
        \item Optimization for classification-specific challenges
    \end{itemize}
\end{enumerate}

\begin{table}[!htb]
\centering
\begin{tabular}{l|cccc}
\toprule
 & H2O & LightAutoML & AutoGluon & FEDOT \\
\midrule
Approach & TAML \& Word2Vec & TAML \& emb. & BERT fine-tuning & TAML \& TF-IDF \\
Scarce train data & $\times$ & Has small data modes & $\times$ & Adaptable \\
Custom search space & Flexible API & HP presets & \checkmark & HP presets \\
Experiments tracking & H2O Flow integration & $\times$ & $\times$ & $\times$ \\
Embedding prompts & $\times$ & $\times$ & $\times$ & $\times$ \\
OOS detection & $\times$ & $\times$ & $\times$ & $\times$ \\
\bottomrule
\end{tabular}
\caption{Comparison of NLP functionality in AutoML frameworks. HP stands for hyperparameters, TAML stands for tabular AutoML.}
\label{tab:frameworks}
\end{table}

\subsection{Structure}

AutoIntent is designed to solve generic text classification problems with comprehensive support for multi-label classification and OOS detection. The framework's architecture is built around a modular design that enables both standalone usage of individual components and end-to-end automated optimization through the AutoML pipeline (fig. \ref{fig:autointent_pipeline}).

\begin{figure}[htbp]
    \centering
    \begin{minipage}{0.45\textwidth}  % First column
        \begin{minted}[
            frame=lines,
            framesep=2mm,
            baselinestretch=1.2,
            bgcolor=white,
            fontsize=\footnotesize,
            linenos,
            numbersep=5pt
        ]{python}
from autointent import Pipeline, Dataset

dataset = Dataset.from_json(path_to_json)
pipeline = Pipeline.from_preset("light")
pipeline.fit(dataset)
pipeline.predict(
  ["show me my latest transactions"]
)
        \end{minted}
        \caption{Example usage of AutoIntent pipeline.}
        \label{fig:autointent_pipeline}
    \end{minipage}%
    \hfill
    \begin{minipage}{0.45\textwidth}  % Second column
        \begin{minted}[
            frame=lines,
            framesep=2mm,
            baselinestretch=1.2,
            bgcolor=white,
            fontsize=\footnotesize,
            linenos,
            numbersep=5pt
        ]{python}
from autointent.modules import KNNScorer

scorer = KNNScorer(
  embedder_config="sergeyzh/rubert-tiny-turbo",
  k=5,
)
scorer.fit(["hello", "how are you?"], [0, 1])
scorer.predict(["hi", "what's up?"])
        \end{minted}
        \caption{AutoIntent's modular approach.}
        \label{fig:knnscorer}
    \end{minipage}
\end{figure}

The framework's key features include:

\begin{itemize}
    \item \textbf{Sklearn-like Interface}: All components follow the familiar sklearn estimator interface \cite{pedregosa2018scikitlearnmachinelearningpython}, making the framework accessible to users familiar with standard machine learning workflows.
    
    \item \textbf{AutoML Pipeline}: The core of AutoIntent is its automated machine learning pipeline, which sequentially optimizes three main components:
    \begin{enumerate}
        \item Embedding Module: Selects and optimizes the most appropriate transformer model from Hugging Face for the given task
        \item Scoring Module: Identifies and tunes the best performing classifier for the task
        \item Decision Module: Optimizes classification thresholds for multi-label scenarios and OOS detection
    \end{enumerate}
    It is implemented with the use of Sentence Transformers library \cite{reimers-2019-sentence-bert}.
    
    \item \textbf{Augmentation Modules}: The framework includes data augmentation capabilities:
    \begin{itemize}
        \item LLM-based augmentation for generating synthetic training examples
        \item Evolutionary approach for transforming existing utterances to expand the training set
        \item Customizable prompts to fine-tune the augmentation process for your needs
    \end{itemize}
    This functionality implemented to be compatible with any LLM via OpenAI completions API \cite{openai_api}.
\end{itemize}

\subsection{Separation of Concerns}

In AutoIntent, a scoring module is defined as a model that takes a text sample as input and outputs a vector of class probabilities (fig. \ref{fig:knnscorer}). This definition is crucial as it establishes a clear separation of concerns in the framework's architecture: scoring modules are responsible for probability estimation, while decision modules take these probabilities as input and produce the final classification decisions (single label, multiple labels, or OOS detection).

This separation of concerns provides several key benefits:

\begin{enumerate}
    \item \textbf{Modularity}: Each component can be developed, tested, and optimized independently. Scoring modules focus solely on learning the relationship between text and class probabilities, while decision modules concentrate on the optimal way to convert these probabilities into final predictions.
    
    \item \textbf{Flexibility}: The same scoring module can be paired with different decision modules to address various scenarios. For example, a single scoring module can be used with:
    \begin{itemize}
        \item A threshold-based decision module for binary classification
        \item An adaptive decision module for multi-label classification
        \item A specialized OOS detection decision module
    \end{itemize}
    
    \item \textbf{Reusability}: Probability estimates from scoring modules can be stored and reused with different decision strategies without recomputing the scores, introducing AutoIntent's efficient experimentation with various decision-making approaches.
    
    \item \textbf{Interpretability}: The separation allows for better understanding of the pipeline's behavior, as probability estimates can be analyzed independently of the final decision-making process.
\end{enumerate}

\subsection{Scoring Module}

AutoIntent offers a diverse set of scoring models, each implementing different classification approaches. A key architectural feature of the framework is that all classifiers (except BERT-based ones) operate on pre-computed embeddings as features. This design choice provides several advantages.

\textbf{Efficiency}: By separating the computationally intensive embedding computation from classification, the framework achieves a fine balance between effectiveness and computational efficiency. The embedding vectors serve as rich, pre-computed features that capture the semantic meaning of the text, while the classifiers themselves are lightweight and fast to train and deploy.
    
\textbf{Hardware Flexibility}: This architecture allows the framework to be deployed on a broad range of machines, even without GPU acceleration. The embedding computation can be done once and cached, while the classification models can run efficiently on CPU-only systems. This is extremely important for hyperparameter optimization.

The available scoring modules include:

\begin{enumerate}
    \item \textbf{Linear classifier}: Utilizes sklearn's \texttt{LogisticRegressionCV} \cite{pedregosa2018scikitlearnmachinelearningpython} for efficient linear classification with built-in cross-validation. The linear model operates directly on the embedding vectors, providing fast training and inference times.
    
    \item \textbf{KNN-based approaches}:
    \begin{itemize}
        \item KNN with FAISS \cite{douze2024faiss} for efficient nearest neighbor search in the embedding space
        \item Rerank: A two-stage approach, essentialy it is a KNN wrapped with cross-encoder reranking and optional cutting off.
        \item MLKNN \cite{zhang2007ml}: A specialized multi-label KNN implementation with Bayesian background, leveraging the embedding space for similarity calculations
    \end{itemize}
    
    \item \textbf{BERT-based classifiers}: full model fine-tuning \cite{devlin2019bertpretrainingdeepbidirectional} and parameter-efficient approaches \cite{hu2021loralowrankadaptationlarge,liu2023gptunderstands}.
    
    \item \textbf{Generic sklearn integration}: Support for any sklearn classifier, allowing for easy extension with additional classification algorithms. All sklearn models operate on the embedding vectors as features, maintaining the efficiency benefits of the architecture.
\end{enumerate}

Each scoring module follows a consistent interface, providing both standard prediction methods and rich output options for detailed analysis and debugging.

\subsection{Decision Module}

The Decision Module handles the final classification decisions, particularly crucial for multi-label scenarios and out-of-sample detection. It processes the probability scores from the Scoring Module to produce the final set of predicted labels. AutoIntent provides several specialized decision strategies:

\begin{enumerate}
    \item \textbf{AdaptiveDecision} \cite{hou2020fewshotlearningmultilabelintent}: A sample-specific thresholding approach that calculates thresholds tailored for each text sample based on the minimum and maximum class scores. This method is specifically designed for multi-label classification scenarios, where it can better handle the varying confidence levels across different samples. However, it does not support out-of-sample detection.
    
    \item \textbf{JinoosDecision} \cite{zhang2020discriminativenearestneighborfewshot}: A specialized approach for handling out-of-sample (OOS) detection, adapted from the DNNc framework. This method finds a universal threshold that balances between in-domain and out-of-domain accuracy, making it particularly suitable for applications where detecting unknown intents is crucial.
    
    \item \textbf{ThresholdDecision}: A simple but effective approach using a fixed threshold specified during model initialization. This method provides a straightforward way to control the classification confidence level and can be useful in scenarios where consistent decision boundaries are required. Also this module is handy to use within AutoML pipeline with optuna samplers, which we will discuss later.
    
    \item \textbf{TunableDecision}: An automated threshold optimization approach that uses Optuna \cite{akiba2019optuna} to find the optimal threshold value by maximizing the F1 score on the given data. This method is particularly useful when the optimal threshold is not known in advance and needs to be determined empirically.
\end{enumerate}

\subsection{Embedding Module}\label{sec:embedding}

AutoIntent leverages the sentence-transformers library, providing access to a wide range of pre-trained transformer models from Hugging Face Hub \cite{wolf2020huggingfacestransformersstateoftheartnatural}. While the embedding module can be used standalone, its primary role is within the AutoML pipeline, where it serves as a single-time selection step of the most appropriate embedding model for the task.

The framework offers three strategies for handling embedding model selection:

\begin{enumerate}
    \item \textbf{Pipeline-level Optimization}: The embedding model is chosen once at the pipeline level, before proceeding to classifier optimization. This approach is more efficient because it avoids the need to refit and reevaluate each scoring module for every candidate embedding model. Users can select from two optimization criteria:
    \begin{itemize}
        \item \textbf{Retrieval-based}: Evaluates embedding models based on their ability to match samples within the same class. Quality is measured using retrieval metrics such as hit rate or NDCG, which can be customized by the user.
        \item \textbf{Classification-based}: Assesses embedding models by their performance with a logistic regression classifier. Uses classification metrics like accuracy or ROC AUC, which can also be customized by the user.
    \end{itemize}
    
    \item \textbf{Scoring-level Optimization}: Users can opt to optimize the embedding model individually for each scoring module. While this approach might yield better results, it is the most time and resource-intensive as it requires refitting and reevaluating each scoring module for every candidate embedding model.
    
    \item \textbf{Fixed Embedding}: Users can specify a default embedding model and skip the optimization process entirely, which is useful when the optimal model is known in advance or when computational resources are limited.
\end{enumerate}

This flexible approach to embedding model selection allows users to balance between optimization quality and computational efficiency based on their specific needs and resource constraints.

\subsection{AutoML Pipeline}

The AutoML Pipeline includes main functionalities: data splitting for hold-out or cross-validation, hyperparameters optimization, models training and saving to disk, classification metrics tracking and logging.

AutoIntent orchestrates the optimization of all components hierarchically (fig. \ref{fig:autointent_pipeline}), with three distinct levels of optimization. At the highest level, the pipeline performs module-level optimization, where it sequentially optimizes the embedding, scoring, and decision modules. Each module builds upon the best model from the previous module's optimization, creating a cohesive pipeline. For instance, the scoring module utilizes features from the best embedding model, while the decision module processes probabilities from the best classifier. This greedy approach effectively prevents combinatorial explosion while maintaining strong performance.

\begin{figure}[!htb]
    \centering
    \includegraphics[width=0.98\linewidth]{pics/autointent_hpo.drawio.pdf}
    \caption{\textbf{(Left)} AutoIntent's three levels of hyperparameter optimization: at the module level, the embedding, scoring, and decision models are optimized sequentially; at the model level, each classification approach is tested against each other to select the best one; at the instance level, hyperparameters for each model is tuned individually with optuna samplers. \textbf{(Right)} Inference pipeline as a result of AutoIntent's hyperparameter optimization.}
    \label{fig:autointent_hpo}
\end{figure}

The second level focuses on model-level optimization, where the pipeline exhaustively evaluates all specified models for each module. This includes various transformer models for the embedding module, different classification approaches (KNN, Logistic Regression, BERT fine-tuning, and sklearn classifiers) for the scoring module, and multiple threshold strategies for the decision module. This comprehensive evaluation ensures that the best model type is selected for each specific task.

At the most granular level, the pipeline performs instance-level optimization. This involves tuning model-specific hyperparameters and optimizing trainable weights where applicable. For example, it tunes parameters such as the number of neighbors in KNN and threshold values, while also training feature weights in Logistic Regression. The optimization process employs Optuna's sampling strategies: random sampling, brute force search, and Tree-structured Parzen Estimators \cite{watanabe2023treestructuredparzenestimatorunderstanding}.

A crucial aspect of the pipeline is its clear distinction between tuning and training. We refer to tuning as the optimization of hyperparameters that control the model's behavior (e.g., number of neighbors in KNN), while training refers to the optimization of model weights learned from the data (e.g., coefficients in Logistic Regression). This distinction is essential for proper model optimization and evaluation. The pipeline implements careful data handling to prevent target leakage and ensure reliable model selection. It maintains separate datasets for different purposes: training data for model weights and parameters, validation data for hyperparameter tuning and model selection, and distinct validation sets for scoring and decision modules to prevent target leakage during holdout scoring. Cross-validation is implemented with the classic approach to an ensemble of stacked models, it utilizes out-of-fold probabilities from the previous layer of models as the data for the next layer of models.

The optimization process is configured through a declarative search space definition, allowing users to specify available modules for each stage, hyperparameter ranges, optimization metrics, and search strategies. The result is a fully optimized classification system that includes the best embedding model for the task, the optimal scoring module with tuned hyperparameters and trained weights, and the appropriate decision module with optimized thresholds. All necessary configuration for inference is provided through intuitive \texttt{save}, \texttt{load}, and \texttt{predict} methods.

\section{Experiments}

\subsection{Scoring Modules}\label{sec:scoring_experiments}

We conducted a comprehensive evaluation of all scoring models across a range of popular intent classification datasets \cite{banking77,minds14,hwu64,snips,massive} using lightweight backbone \cite{he2021debertav3} and embedding models \cite{reimers-2019-sentence-bert}. The experimental setup employed holdout validation and hyperparameter optimization using the TPE sampler \cite{watanabe2023treestructuredparzenestimatorunderstanding} with a maximum of 20 trials per model (see complete search space specifications in Appendix \ref{app:search_spaces}).

The results, presented in Table \ref{tab:scoring_light}, reveal several key findings. Logistic regression over transformer embeddings (<<linear>> in the table) demonstrates superior performance in terms of both average accuracy and consistency across datasets, achieving the highest score in three out of five datasets. BERT-based methods exhibit significant variability in performance: parameter-efficient approaches (P-tuning and LoRA) show notably lower performance, full BERT fine-tuning achieves competitive results in one dataset but underperforms in others. This suggests potential challenges in hyperparameter optimization for fine-tune-based methods. Zero-shot-based methods (linear, KNN, random forest) consistently outperform transformer-based approaches, with all three achieving average accuracy above 90\%.

\begin{table}[h]
\begin{tabular}{l|rrrrr|r}
\toprule
model name & banking77 & minds14 & hwu64 & snips & massive & average \\
\midrule
ptuning & 4.63 & 11.63 & 3.65 & 66.49 & 8.57 & 18.99 \\
lora & 20.35 & 12.40 & 22.40 & 95.09 & 41.95 & 38.44 \\
bert & 64.14 & 69.77 & 73.40 & \textbf{98.52} & 76.76 & 76.52 \\
rerank & 89.04 & \textbf{97.67} & 84.45 & 96.59 & 81.68 & 89.89 \\
random forest & 89.81 & 98.45 & 86.98 & 95.16 & 79.73 & 90.03 \\
knn & 89.74 & \textbf{97.67} & 85.42 & 96.10 & 81.65 & 90.12 \\
linear & \textbf{90.51} & \textbf{97.67} & \textbf{89.17} & 97.45 & \textbf{84.70} & \textbf{91.90} \\
\bottomrule
\end{tabular}
\centering
\caption{Accuracy of scoring modules with lightweight backbone. BERT-based methods (bert, lora, ptuning) utilize \texttt{prajjwal1/bert-tiny}, \cite{bhargava2021generalization,tinybert} as a backbone. The rest of methods use embeddings from \texttt{sentence-transformers/all-MiniLM-L6-v2} \cite{reimers-2019-sentence-bert}. Each metric value is obtained via averaging three runs with different seeds with random validation split formed from fixed training data.}
\label{tab:scoring_light}
\end{table}

To investigate the impact of model capacity on performance, we repeated the experiment using heavier transformer models. While the overall ranking by average accuracy remained consistent with the lightweight experiment, the distribution of best-performing models across datasets shifted. The dominance of logistic regression became less pronounced, with KNN and random forest achieving comparable performance in terms of dataset-specific wins. BERT-based methods continued to underperform relative to feature-based approaches, though full BERT fine-tuning showed improved results in some datasets. This suggests that increased model capacity alone may not address the optimization challenges of transformer-based methods. The competitive performance of multiple feature-based methods justifies the inclusion of diverse scoring modules in the library, particularly for applications where identifying the optimal model is critical.

\begin{table}[h]
\centering
\label{tab:scoring_heavy}
\begin{tabular}{l|rrrrr|r}
\toprule
model name & banking77 & minds14 & hwu64 & snips & massive & average \\
\midrule
ptuning & 5.00 & 14.73 & 9.64 & 91.67 & 25.87 & 29.38 \\
lora & 13.82 & 9.30 & 38.58 & 99.13 & 52.06 & 42.58 \\
bert & 85.61 & 48.06 & 88.58 & \textbf{99.16} & 86.36 & 81.56 \\
rerank & 90.28 & 96.12 & 87.24 & 97.35 & 83.87 & 90.97 \\
random forest & 91.54 & \textbf{98.45} & 90.22 & 96.64 & 85.08 & 92.39 \\
knn & \textbf{92.14} & \textbf{98.45} & 89.43 & 96.97 & 85.43 & 92.48 \\
linear & 91.81 & 96.90 & \textbf{91.78} & 98.06 & \textbf{87.58} & \textbf{93.23} \\
\bottomrule
\end{tabular}
\caption{Accuracy of scoring modules with lightweight backbone. BERT-based methods (bert, lora, ptuning) utilize \texttt{microsoft/deberta-v3-small}, \cite{he2021debertav3} as a backbone. The rest of methods use embeddings from \texttt{mixedbread-ai/mxbai-embed-large-v1} \cite{emb2024mxbai,li2023angle}. Each metric value is obtained via averaging three runs with different seeds with random validation split formed from fixed training data.}
\end{table}

\subsection{Computational Efficiency}

To quantify the computational requirements of different scoring modules, we conducted a comprehensive analysis using the Code Carbon library \cite{benoit_courty_2024_11171501}. This analysis measured various aspects of computational resource consumption for a single trial (training and evaluation of a single model configuration). The results, presented in Table \ref{tab:computational}, reveal significant variations in resource usage across different approaches.

\begin{table}[h]
\centering
\begin{tabular}{l|ccccccc}
\toprule
model & emissions & runtime & energy & gpu & cpu & ram & rate \\
\midrule
bert & 1.382 & 103.911 & 3.133 & 2.198 & 0.774 & 1.615e-01 & 0.014 \\
ptuning & 1.118 & 83.455 & 2.535 & 1.785 & 0.620 & 1.295e-01 & 0.014 \\
lora & 0.863 & 65.157 & 1.957 & 1.372 & 0.484 & 1.009e-01 & 0.013 \\
linear & 0.428 & 73.393 & 0.971 & 0.312 & 0.545 & 1.138e-01 & 0.006 \\
rerank & 0.270 & 29.040 & 0.613 & 0.355 & 0.213 & 4.436e-02 & 0.010 \\
dnnc & 0.122 & 10.000 & 0.276 & 0.192 & 0.070 & 1.455e-02 & 0.013 \\
rand forest & 0.073 & 11.367 & 0.166 & 0.074 & 0.080 & 1.664e-02 & 0.007 \\
knn & 0.009 & 1.281 & 0.019 & 0.014 & 0.004 & 9.044e-04 & 0.012 \\
\midrule
units & grams & sec & Wh & Wh & Wh & Wh & grams/sec \\
\bottomrule
\end{tabular}
\caption{Computational resource consumption for different scoring modules. The experiments are conducted on banking77 dataset with \texttt{mixedbread-ai/mxbai-embed-large-v1} \cite{emb2024mxbai,li2023angle}, system with AMD Ryzen 7 5800H, NVIDIA RTX 3060 Laptop. Median values of 10 trials are displayed. Embeddings were pre-computed.}
\label{tab:computational}
\end{table}

The analysis revealed significant variations in resource efficiency across different scoring methods. KNN-based methods demonstrated exceptional efficiency, with minimal emissions and runtime, making them particularly suitable for resource-constrained environments. Logistic regression showed moderate resource consumption while maintaining high performance, representing a balanced trade-off between computational cost and effectiveness. In contrast, BERT-based methods exhibited the highest resource requirements, necessitating substantial computational infrastructure. These findings provide valuable insights for deployment scenarios with varying resource constraints and for building AutoIntent's presets.

\subsection{Evolutionary Augmentations}\label{sec:augmentations}

To evaluate the effectiveness of LLM-based data augmentation in real-world scenarios with limited training data, we conducted our experiments on randomly sampled 10-shot versions of the datasets. This setup simulates common practical situations where end users need to expand their small training splits. We conducted a comprehensive correlation analysis between the number of synthetic samples added and the final classification accuracy. We tested four state-of-the-art LLMs (\texttt{DeepSeek-V3-0324} \cite{deepseekai2025deepseekv3technicalreport}, \texttt{GPT-4o-mini} \cite{gpt4o_blog}, \texttt{Qwen2.5-7B} \cite{qwen2025qwen25technicalreport}, and \texttt{Llama-3.1-8B} \cite{llama31_blog}) across multiple datasets. The correlation analysis (tab. \ref{tab:correlation}, \ref{tab:dataset_correlation}) revealed several key findings.

\textbf{Model-specific Patterns}. Surprisingly, the size and type of the LLM (proprietary vs. open-source) showed no clear correlation with augmentation effectiveness. Both smaller models like Llama-3.1-8B and larger models like DeepSeek-V3-0324 demonstrated similar patterns of effectiveness.

These results suggest that the effectiveness of augmentation doesn't vary significantly across different LLMs, but with GPT-4o-mini showing the most consistent positive impact. English datasets generally benefit more from augmentation than their Russian counterparts, likely due to the English-oriented nature of the prompts we used for augmentation. OOS detection performance tends to decrease with increased augmentation, particularly for certain models. This can be attributed to the combination of factors: the Russian language of the dataset, its high granularity (150 classes), and the generic nature of synthetic samples that fail to capture the fine-grained class boundaries. The banking77 dataset shows the most consistent improvement with augmentation across all models

\begin{table}[h]
\centering
\begin{minipage}{0.6\textwidth}
\centering
\begin{tabular}{l|rrr}
\toprule
                           Model & English & Russian  &  OOS \\
                           \midrule 
                DeepSeek-V3-0324 &  $0.532$  &     -  & $-0.777$ \\
          gpt-4o-mini-2024-07-18 &  $0.500$  &     -  & $0.884$ \\
         Qwen2.5-7B-Instruct-AWQ &      -  &     -  & $-0.613$ \\
Meta-Llama-3.1-8B-Instruct-Turbo &  $0.568$  &     -  & $-0.932$ \\
\bottomrule
\end{tabular}
\caption{Correlation Analysis (Spearman r values, non-significant correlations are replaced with dash). The support for calculating each value in <<English>> and <<Russian>> columns is 22 points both, for <<OOS>> it is 11. English datasets: banking77 \cite{banking77}, snips \cite{snips}; Russian datasets banking77, snips (manually translated by us); OOS datasets: clinc150 \cite{clinc150} translated into Russian. See Appendix \ref{app:detailed_aug_metrics} for detailed metrics report.}
\label{tab:correlation}
\end{minipage}
\hfill
\begin{minipage}{0.35\textwidth}
\centering
\begin{tabular}{l|r}
\toprule
     Dataset & Correlation \\
    \midrule
   banking77 \cite{banking77} &      $0.710$ \\
   snips \cite{snips} &      $0.459$ \\
banking77\_ru &          - \\
    snips\_ru &          - \\   
 clinc150\_ru &          - \\
\bottomrule
\end{tabular}
\caption{Dataset-specific correlations across all models (Spearman r values, non-significant correlations are replaced with dash). The support for calculating each value is 44.}
\label{tab:dataset_correlation}
\end{minipage}
\end{table}

\begin{wrapfigure}[22]{r}{0.45\textwidth}
    \centering
    \includegraphics[width=1\linewidth]{pics/mean_improvement.pdf}
    \caption{Mean accuracy improvement compared to number of synthetic examples added to 10-shot versions of datasets: banking77 \cite{banking77}, snips \cite{snips} and their Russian translations. The support for calculating each value in this figure is 16 points.}
    \label{fig:augmentation_analysis}
\end{wrapfigure}

To further investigate the impact of augmentation on model performance, we conducted a detailed statistical analysis comparing the accuracy before and after augmentation for each augmentation level (1-10 synthetic samples added to train set). We performed paired t-tests to assess the statistical significance of the improvements and calculated effect sizes using Cohen's d to quantify the magnitude of the changes (fig. \ref{fig:augmentation_analysis}).

The analysis revealed that all augmentation levels resulted in statistically significant improvements ($p \leq 0.002$) compared to the baseline (no augmentation). The effect sizes (Cohen's d) ranged from $0.35$ to $0.62$, indicating moderate to large practical significance. The mean improvement in accuracy showed a generally increasing trend with the number of augmentations, ranging from 7.1\% to 11.7\% improvement over the baseline.

These results demonstrate the following. First, any amount of augmentation (1-10 samples) leads to statistically significant improvements in model performance. Second, the relationship between the number of augmentations and performance improvement is not strictly linear, with quite early saturation. Third, the improvements are not only statistically significant but also practically meaningful, as indicated by the moderate to large effect sizes.

This analysis provides strong empirical support for the effectiveness of LLM-based augmentation in improving classification performance, while also suggesting optimal ranges for the number of synthetic samples to generate.

\subsection{Baselines}\label{sec:baselines_exps}

To evaluate AutoIntent's performance, we conducted a comparative analysis against several open-source NLP AutoML frameworks. The baseline frameworks were selected based on their availability and relevance to text classification tasks:

\begin{enumerate}
    \item \textbf{H2O} \cite{ledell2020h2o}: Implements tabular AutoML methods over Word2Vec embeddings;
    \item \textbf{LightAutoML (LAMA)} \cite{vakhrushev2022lightautomlautomlsolutionlarge}: Performs AutoML over TF-IDF features with optional support for three predefined transformer models;
    \item \textbf{FEDOT} \cite{nikitin2021automated}: Utilizes TF-IDF based feature engineering;
    \item \textbf{AutoGluon} \cite{tang2024autogluonmultimodalautommsuperchargingmultimodal}: Employs a fixed training recipe with DeBERTa-v3 \cite{he2021debertav3} using PyTorch Lightning.
\end{enumerate}

The evaluation was conducted across five standard intent classification datasets, with results presented in Table \ref{tab:baselines}. AutoIntent's configuration matched the heavy backbone experiment, but with consistent test data across all frameworks.

\begin{table}[h]
\centering
\begin{tabular}{l|rrrrr|r}
\toprule
framework & banking77 & hwu64 & massive & minds14 & snips & avg \\
\midrule
autogluon & \textbf{93.28} & \textbf{91.17} & \textbf{88.92} & \textbf{97.22} & \textbf{99.07} & \textbf{93.13} \\
autointent & 92.86 & 90.83 & 87.13 & 95.68 & 98.19 & 92.94 \\
h2o & 75.32 & 77.32 & 75.30 & 76.85 & 98.36 & 80.63 \\
fedot & 1.30 & 1.77 & 7.28 & 12.04 & 15.14 & 7.51 \\
lama & 1.30 & 1.77 & 7.04 & 8.33 & 14.50 & 6.59 \\
\bottomrule
\end{tabular}
\caption{Performance comparison across different AutoML frameworks.}
\label{tab:baselines}
\end{table}

The results reveal several key findings:

\begin{enumerate}
    \item \textbf{Framework Performance}:
    \begin{itemize}
        \item AutoGluon and AutoIntent demonstrate comparable performance
        \item H2O shows moderate performance
        \item FEDOT and LAMA fail to achieve meaningful results across most datasets
    \end{itemize}
    
    \item \textbf{Framework Limitations}:
    \begin{itemize}
        \item \textbf{Hyperparameter Optimization}: AutoGluon uses a fixed training recipe without hyperparameter tuning (note: attempted HPO presets encountered unexpected bugs)
        \item \textbf{Feature Engineering}: H2O and FEDOT lack support for transformer-based features
        \item \textbf{Model Flexibility}: LAMA's transformer support is limited to three predefined models
        \item \textbf{Inference Efficiency}: AutoGluon outputs fixed inference-time models, while AutoIntent can select lighter models when performance is comparable
        \item \textbf{Deployment Features}: Most of the baseline frameworks lack mature deployment capabilities, whereas AutoIntent provides straightforward model serialization
    \end{itemize}
\end{enumerate}

The comparative analysis highlights AutoIntent's advantages in terms of flexibility, deployment capabilities, and adaptive model selection, while maintaining competitive performance with state-of-the-art approaches.

\subsection{OOS Detection}\label{sec:oos_detection}

To evaluate AutoIntent's out-of-scope (OOS) detection capabilities, we conducted a comparative analysis on the CLINC150 dataset \cite{clinc150}, which is specifically designed for intent classification with OOS detection. Since all the baseline frameworks lack native OOS detection support, we adapted their configurations by treating OOS as an additional class. The results are presented in Table \ref{tab:oos}.

\begin{table}[h]
\centering
\begin{tabular}{l|ccc}
\toprule
framework & in domain acc & OOS precision & OOS recall \\
\midrule
autointent & \textbf{96.13} & 91.92 & \textbf{65.93} \\
autogluon \cite{tang2024autogluonmultimodalautommsuperchargingmultimodal} & 95.76 & \textbf{98.47} & 32.20 \\
h2o \cite{ledell2020h2o} & 85.22 & 82.57 & 27.00 \\
fedot \cite{nikitin2021automated} & 0.00 & 18.18 & 100.0 \\
lama \cite{vakhrushev2022lightautomlautomlsolutionlarge} & 0.00 & 18.18 & 100.0 \\
\bottomrule
\end{tabular}
\caption{Performance comparison on OOS detection task.}
\label{tab:oos}
\end{table}

These results demonstrate the superiority of AutoIntent in solving the OOS detection problem due to the implementation of the threshold tuning stage.

\subsection{Embedding Module Effectiveness}

We evaluated the effectiveness of AutoIntent's feature of selecting embedding model via optimizing embedding module as it was described in Section \ref{sec:embedding}. Figures \ref{tab:encoder_ranking},\ref{fig:encoder_ranking} contain results for two scenarios of evaluation to select best embedding model: the first is to run the whole AutoML pipeline and measure final accuracy, the second is to use our embedding module to optimize NDCG of retrieval quality. These results show the imperfection of the approximate embedding model selection, which is expected as it is done on a subset of a data and without full pipeline training. However, the general ranking is quite close and the best model (\texttt{stella\_en\_400M\_v5}) is chosen correctly, so the optimizing embedding module reaches good balance between the computation resources spent and the correctness of the final result obtained.

\begin{figure}[h]
\centering
\begin{minipage}[t]{0.5\textwidth}
\vspace{0pt}
\begin{tabular}{l|cc}
\toprule
\textbf{Model} & \textbf{Accuracy} & \textbf{NDCG} \\
\midrule
stella\_en\_400M\_v5 & 94.28 & 93.83 \\
multilingual-e5-large-instruct & 93.65 & 92.97 \\
GIST-large-Embedding-v0 & 93.51 & 93.32 \\
UAE-Large-V1 & 92.89 & 93.25 \\
bge-m3 & 92.69 & 92.49 \\
multilingual-e5-large & 91.41 & 92.45 \\
LaBSE & 90.47 & 89.51 \\
KaLM-embedding-* & 89.65 & 92.88 \\
nomic-embed-text-v1.5 & 87.24 & 89.63 \\
deberta-v3-small & 81.15 & 67.20 \\
deberta-v3-large & 75.39 & 59.59 \\
deberta-v3-base & 75.00 & 59.28 \\
\bottomrule
\end{tabular}
\caption{Embedding models performance averaged over hwu64 \cite{hwu64}, massive \cite{massive}, minds14 \cite{minds14}, snips \cite{snips}. $^*$multilingual-mini-instruct-v1.5. Hugging Face Hub organizations in the corresponding order: NovaSearch/, intfloat/, avsolatorio/, WhereIsAI/, BAAI/, intfloat/, sentence-transformers/, HIT-TMG/, nomic-ai/, microsoft/, microsoft/, microsoft/.}
\label{tab:encoder_ranking}
\end{minipage}
\hfill
\begin{minipage}[t]{0.45\textwidth}
\vspace{0pt}
\centering
\includegraphics[width=1\linewidth]{pics/encoders_ranking.pdf}
\caption{Encoders ranking: (\textbf{Left}) precise ranking obtained via training full AutoML pipeline with only this model, (\textbf{Right}) approximate ranking based on retrieval quality (NDCG).}
\label{fig:encoder_ranking}
\end{minipage}
\end{figure}

\section{AutoIntent Cookbook}

Based on our extensive experimental analysis, we present a comprehensive guide for selecting and configuring AutoIntent components for different use cases. To facilitate the adoption of these recommendations, we have implemented them as configurable presets in our library.

\subsection{Resource-aware Configuration Presets}

We have identified three distinct configuration presets that balance performance and computational requirements:

\begin{enumerate}
    \item \textbf{Lightweight Configuration}:
    \begin{itemize}
        \item \textbf{Components}: KNN-based methods
        \item \textbf{Use Case}: 
        \begin{itemize}
            \item Scenarios requiring minimal resource consumption
            \item Few-shot learning scenarios with limited training data
        \end{itemize}
        \item \textbf{Advantages}:
        \begin{itemize}
            \item Fast inference time
            \item Minimal computational requirements for training
            \item Particularly effective in few-shot learning scenarios
        \end{itemize}
        \item \textbf{Performance}: 
        \begin{itemize}
            \item Achieves competitive accuracy with significantly reduced resource usage
            \item Demonstrates strong performance on small datasets, often outperforming more complex models in few-shot settings
            \item Maintains consistent performance across different data sizes
        \end{itemize}
    \end{itemize}
    
    \item \textbf{Moderate Configuration}:
    \begin{itemize}
        \item \textbf{Components}: Linear scorer, cross-encoder based methods, and all lightweight methods
        \item \textbf{Use Case}: Applications requiring consistent, high-quality performance
        \item \textbf{Advantages}:
        \begin{itemize}
            \item Balanced performance and resource usage
            \item Robust across different datasets
            \item Suitable for production environments
        \end{itemize}
        \item \textbf{Performance}: Provides reliable accuracy with moderate computational overhead
    \end{itemize}
    
    \item \textbf{Heavy Configuration}:
    \begin{itemize}
        \item \textbf{Components}: BERT-based models, along with all moderate and lightweight methods
        \item \textbf{Use Case}: High-performance scenarios with access to professional-grade GPUs
        \item \textbf{Requirements}:
        \begin{itemize}
            \item High-performance GPU with substantial VRAM
            \item Extended training time
            \item Higher computational budget
            \item Large dataset
        \end{itemize}
        \item \textbf{Performance}: Potential for superior accuracy with heavy backbones and longer training
    \end{itemize}
\end{enumerate}

\subsection{Selection Guidelines}

When choosing a configuration preset, consider the following factors:

\begin{enumerate}
    \item \textbf{Resource Constraints}:
    \begin{itemize}
        \item Available computational resources
        \item Inference time requirements
        \item Memory limitations
    \end{itemize}
    
    \item \textbf{Performance Requirements}:
    \begin{itemize}
        \item Required accuracy
        \item Latency constraints
        \item Scalability needs
    \end{itemize}
    
    \item \textbf{Dataset Characteristics}:
    \begin{itemize}
        \item Size of the training data
        \item Complexity of the classification task
        \item Language and domain specifics
    \end{itemize}
\end{enumerate}

\subsection{Implementation Notes}

Each preset is implemented as a configurable template in the library, allowing users to:
\begin{itemize}
    \item Start with recommended configurations
    \item Customize parameters based on specific needs
    \item Gradually scale up from lightweight to heavy configurations as requirements evolve
\end{itemize}

This cookbook approach enables users to quickly deploy effective solutions while maintaining the flexibility to adapt to changing requirements and resource constraints.

\section{Future Research}

Our experimental analysis and framework development have revealed several promising directions for future research:

\begin{enumerate}
    \item \textbf{Framework Evaluation Extensions}:
    \begin{itemize}
        \item Comprehensive evaluation of module data separation strategies
        \item In-depth analysis of multi-label classification performance
        \item Systematic assessment of out-of-scope (OOS) detection capabilities
        \item Investigation of cross-lingual transfer capabilities
    \end{itemize}
    
    \item \textbf{Augmentation Method Improvements}:
    \begin{itemize}
        \item Development of more appropriate prompt engineering techniques
        \item Integration of domain-specific knowledge into augmentation
        \item Exploration of hybrid approaches combining LLM-based and traditional augmentation methods
        \item Investigation of augmentation quality control mechanisms
        \item Research into language-specific augmentation strategies
        \item Research on RL-style augmentations with actor and critic
    \end{itemize}
    
    \item \textbf{BERT-based Methods Enhancement}:
    \begin{itemize}
        \item Evaluation of transformer models on high-performance GPUs (e.g., A100)
        \item Investigation of potential overfitting issues in current implementations
    \end{itemize}
    
    \item \textbf{Pipeline Optimization}:
    \begin{itemize}
        \item Development of more optimized module selection strategies
        \item Investigation of distributed HPO approaches
    \end{itemize}
    
    \item \textbf{Deployment and Scalability}:
    \begin{itemize}
        \item Research into efficient model serving strategies
        \item Investigation of distributed HPO approaches
        \item Development of automated scaling mechanisms
        \item Exploration of edge deployment capabilities
    \end{itemize}
\end{enumerate}

These research directions aim to address current limitations and expand the framework's capabilities, making it more robust and adaptable to various real-world scenarios.

\section{Limitations}

\section{Conclusion}

In this report, we presented AutoIntent, an end-to-end AutoML framework designed specifically for text classification. By integrating transformer-based embeddings, a variety of scoring and decision modules, and robust hyperparameter optimization strategies, AutoIntent addresses key limitations of existing NLP AutoML solutions. The framework emphasizes modularity, extensibility, and ease of deployment, supporting use cases from lightweight inference scenarios to high-performance model training.

Our extensive experimental evaluation demonstrates that AutoIntent achieves competitive or superior classification performance across a range of benchmark datasets, while also offering enhanced support for multi-label classification and out-of-scope (OOS) detection. Moreover, its compatibility with large language models for data augmentation and its resource-aware configuration presets make it adaptable to a wide spectrum of real-world scenarios.

AutoIntent combines practical usability with advanced optimization strategies, bridging the gap between research-grade experimentation and production-ready deployment. With its scalable architecture and support for continuous improvement, it provides a solid foundation for future work in automated text classification and conversational AI.


\bibliographystyle{unsrt}  
\bibliography{references} 

\begin{appendices}

\section{Search Spaces}\label{app:search_spaces}

The most clear way to set custom search space within AutoIntent is with yaml configuration. In the sections \ref{sec:scoring_experiments}, \ref{sec:baselines_exps} we used the following search space:
\begin{verbatim}
- node_type: scoring
  target_metric: scoring_accuracy
  search_space:
    - module_name: knn
      n_trials: 20
      k:
        low: 1
        high: 20
      weights: [uniform, distance, closest]
    - module_name: linear
    - module_name: dnnc
      k:
        low: 1
        high: 5
    - module_name: rerank
      k: [30]
      weights: [uniform, distance]
      m: [ 1, 2, 3, 4, 5 ]
      use_crosencoder_scores: [True, False]
    - module_name: sklearn
      n_trials: 20
      clf_name: [RandomForestClassifier]
      n_estimators: [200, 300, 500]
      max_depth: [100, 150, 200]
      max_features: [sqrt, log2, null]
      n_jobs: [8]
    - module_name: bert
      n_trials: 20
      num_train_epochs: [3]
      batch_size: [8, 16, 32]
      learning_rate: [5.0e-5, 1.0e-4, 5.0e-4, 1.0e-3]
    - module_name: lora
      n_trials: 20
      num_train_epochs: [3]
      batch_size: [8, 16, 32]
      learning_rate: [5.0e-5, 1.0e-4, 5.0e-4, 1.0e-3]
      lora_alpha: [16, 32, 64]
      lora_dropout: [0.1, 0.2, 0.3]
      r: [8, 16, 32, 64]
    - module_name: ptuning
      n_trials: 20
      num_train_epochs: [3]
      batch_size: [8, 16, 32]
      learning_rate: [5.0e-5, 1.0e-4, 5.0e-4, 1.0e-3]
      num_virtual_tokens: [10, 30, 50]
      encoder_dropout: [0.1, 0.2, 0.3]
      encoder_hidden_size: [128, 256, 512]
      encoder_reparameterization_type: [MLP, LSTM]
      encoder_num_layers: [1, 2, 3]
- node_type: decision
  target_metric: decision_accuracy
  search_space:
    - module_name: argmax
\end{verbatim}

In the section \ref{sec:oos_detection} we used the same search space for scoring module and the following for decision module:
\begin{verbatim}
- node_type: decision
  target_metric: decision_accuracy
  search_space:
    - module_name: argmax
    - module_name: threshold
      thresh:
        low: 0.1
        high: 0.9
        step: 0.1
    - module_name: jinoos
\end{verbatim}


\section{Augmentation Experiment} \label{app:detailed_aug_metrics}

We report metrics for each dataset-LLM pair from section \ref{sec:augmentations} in figure \ref{fig:detailed_aug_metrics}.

\begin{figure}
    \centering
    \includegraphics[width=0.95\linewidth]{pics/augmented_models_performance.pdf}
    \caption{Detailed metrics for each dataset-LLM pair.}
    \label{fig:detailed_aug_metrics}
\end{figure}

\end{appendices}

\end{document} 